\section*{Bab \ref{ch:b1:ecosystem-dynamics} --- Dinamika Ekosistem}

\begin{solution}{exo:b1:ecosystem-dynamics:1}
Primer: pada alas baru tanpa tanah (sebuah moraine, sebuah aliran
lava); sekunder: tempat tanah dan bank bijinya tetap ada (sebuah
ladang terbengkalai, sebuah hutan terbakar). Yang sekunder lebih
cepat, sebesar faktor sepuluh atau lebih, sebab tanahnya beserta hara
dan bijinya tidak harus dibuat.
\end{solution}

\begin{solution}{exo:b1:ecosystem-dynamics:2}
Fasilitasi, inhibisi, toleransi. Pada lava, fasilitasi: perintisnya
(\omterm{def:b1:species-interactions:symbiosis}{lumut kerak}, lumut, pemancang nitrogen) harus membangun tanah sebelum
apa pun yang lain dapat berakar.
\end{solution}

\begin{solution}{exo:b1:ecosystem-dynamics:3}
\omterm{def:b1:ecosystem-dynamics:disturbance}{Ketahanan}: tinggal tak berubah sepanjang sebuah \omterm{def:b1:ecosystem-dynamics:disturbance}{gangguan}; \omterm{def:b1:ecosystem-dynamics:disturbance}{kelentingan}:
kembali dengan cepat sesudahnya. Sebuah terumbu tidak tahan terhadap
pemanasan maupun lenting; sebuah padang rumput memiliki \omterm{def:b1:ecosystem-dynamics:disturbance}{ketahanan} yang
rendah terhadap sebuah kebakaran dan \omterm{def:b1:ecosystem-dynamics:disturbance}{kelentingan} yang tinggi.
\end{solution}

\begin{solution}{exo:b1:ecosystem-dynamics:4}
$16^{0.25} = 2$: \num{24} spesies; $625^{0.25} = 5$: \num{60} spesies.
\end{solution}

\begin{solution}{exo:b1:ecosystem-dynamics:5}
Endapannya hampir tak bernitrogen; cemaranya tak dapat memancangkannya
lalu tumbuh hanya tempat ia dipasok. Aldernya, yang memancangkan
nitrogen lalu menaikkan kandungan tanahnya menjadi ratusan kilogram
tiap hektar, karena itu adalah yang membuat tapaknya pantas bagi
cemaranya: fasilitasi.
\end{solution}

\begin{solution}{exo:b1:ecosystem-dynamics:6}
$S^* = I_0 P/(I_0 + eP) = 1600/(8 + 10) = 89$ spesies; pergantiannya
$eS^* = 4.4$ spesies setahun. Dengan $I_0 = 4$: $800/14 = 57$ spesies,
pergantiannya \num{2.9}.
\end{solution}

\begin{solution}{exo:b1:ecosystem-dynamics:7}
\omterm{def:b1:ecosystem-dynamics:decomposition}{Penguraiannya} terus memineralkan \omterm{def:b1:ecosystem-dynamics:decomposition}{seresah} dan humusnya menjadi amonium,
dan penitrifikasinya mengubahnya menjadi nitrat, sebuah anion larut
yang tak dipegang tanahnya; tanpa \omterm{def:b1:flowering-plant-organization:organs}{akar} yang menyerapnya, ia tercuci
keluar. Pertumbuhan kembalinya memulihkan serapannya: tumbuhan barunya
mengambil nitratnya secepat ia terbentuk, dan sungainya mengalir
jernih lagi.
\end{solution}

\begin{solution}{exo:b1:ecosystem-dynamics:8}
$k = \ln 2/1 = \qty{0.69}{yr^{-1}}$; sesudah tiga tahun $e^{-2.08} =
\qty{12.5}{\%}$ tersisa. Konifer: $k = -\ln 0.85 = \qty{0.16}{yr^{-1}}$,
waktu paruhnya \num{4.3} tahun, \qty{61}{\%} tersisa sesudah tiga tahun
--- \omterm{def:b1:ecosystem-dynamics:decomposition}{seresahnya} bertimbun.
\end{solution}

\begin{solution}{exo:b1:ecosystem-dynamics:9}
Tanpa penggembalaan: rerumputan tinggi yang bersaing menaungi sisanya,
sedikit spesies. Penggembalaan sedang: yang mendominasi dipangkas,
cahaya mencapai tanahnya, herba kecil berdampingan dengan rerumputan:
keanekaragaman maksimum. Penggembalaan berat: hanya spesies yang
merunduk, tak enak dimakan atau tumbuh cepat yang sintas dari injakan
dan pangkasannya: sedikit spesies lagi.
\end{solution}

\begin{solution}{exo:b1:ecosystem-dynamics:10}
Aslinya: $c\,1000^{0.25} = 5.6c$; satu fragmennya $10^{0.25}c = 1.78c$,
\qty{32}{\%} aslinya. Sepuluh fragmen terasing memegang, seandainya
daftarnya seluruhnya berbeda, $17.8c$ --- lebih banyak daripada
aslinya, yang mustahil: fragmennya berbagi sebagian besar spesiesnya,
sehingga jumlahnya terhitung berlebih. Total sejatinya terletak antara
$1.78c$ (daftar yang serupa) dan sekitar $5.6c$ (aslinya), dikurangi
apa yang disingkirkan keterasingan dan tepinya; spesies yang pertama
hilang adalah yang memerlukan lebih dari \qty{10}{km^2} masing-masing.
\end{solution}

\begin{solution}{exo:b1:ecosystem-dynamics:11}
Sebuah pemulihan yang tepat akan menyarankan sebuah \omterm{def:b1:ecosystem-organization:ecosystem}{komunitas} tetap
yang ditentukan keadaan pulaunya; pemulihan cacahnya dengan sebuah
daftar yang berbeda, dan sebuah daftar yang terus berubah, adalah apa
yang diramalkan sebuah neraca antara imigrasi dan kepunahan yang acak
--- cacahnya adalah sebuah sifat luas dan jaraknya, jati dirinya
adalah kebetulan.
\end{solution}

\begin{solution}{exo:b1:ecosystem-dynamics:12}
Tidak tertutup: ia memasukkan hujan, debu, nitrogen dari udaranya
(lewat fiksasi dan pengendapan) dan ion dari batu yang melapuk, lalu
kehilangan air, ion terlarut dan gas lewat sungai dan udaranya; tetapi
masukan dan keluarannya kecil dibandingkan daur dalamnya dari \omterm{def:b1:ecosystem-dynamics:decomposition}{seresah}
ke tanah ke \omterm{def:b1:flowering-plant-organization:organs}{akar} ke \omterm{def:b1:flowering-plant-organization:organs}{daun}, sehingga hutannya tampak tertutup pada
ketelitian anggaran setahun. Penebangannya memperlihatkan dari mana
ketertutupan yang tampak itu datang: serapan tumbuhan hidupnya; tanpa
serapan itu, tanah yang sama dan \omterm{def:b1:ecosystem-organization:trophic}{pengurai} yang sama membuat lembahnya
membocorkan haranya dalam satu musim.
\end{solution}

\begin{solution}{pb:b1:ecosystem-dynamics:1}
\textbf{1.} 10--40 tahun: $140/30 = \qty{4.7}{kg/ha}$ setahun; 40--80:
$150/40 = \num{3.75}$; 80--200: $20/120 = \num{0.17}$. Paling cepat
pada tahap aldernya.
\textbf{2.} Fasilitasi, bagi alder ke cemaranya: nitrogen yang
ditambahkan aldernya adalah apa yang diperlukan cemaranya.
\textbf{3.} Hutannya telah mencapai sebuah keadaan mantap: masukannya
(fiksasi, pengendapan) mengimbangi kehilangannya (pencucian,
denitrifikasi, kayu yang diambil) dan nitrogennya berdaur di dalamnya
alih-alih bertimbun.
\textbf{4.} \qty{150}{kg/ha} tertahan sepanjang 40 tahun lawan
\num{2000} yang dipancangkan: \qty{7.5}{\%}; sisanya dipegang dalam
biomassa dan \omterm{def:b1:ecosystem-dynamics:decomposition}{seresah} aldernya sendiri, tercuci, atau terdenitrifikasi.
\textbf{5.} Tanamlah semai cemara pada endapan berumur 20 tahun dengan
dan tanpa alder serta dengan dan tanpa pupuk nitrogen: bila cemaranya
tumbuh tanpa alder ketika dipupuk, dan tidak tanpa pupuk, pengaruh
aldernya adalah nitrogennya (fasilitasi); bila cemaranya tumbuh sama
saja pada seluruh petaknya dan sekadar lebih lambat daripada aldernya,
toleransi.
\textbf{6.} Ia menganggap tiap tapaknya bermula serupa lalu mengikuti
jalan yang sama, yang mungkin dipatahkan iklim, kedatangan yang
kebetulan dan endapan gletsernya yang beragam; hanya sebuah petak
tetap yang membuktikan urutannya pada satu tempat.
\textbf{7.} Sebesar \qty{330}{mm}, \qty{40}{\%}: pohonnya tak lagi
mentranspirasikan bagiannya dari hujannya.
\textbf{8.} Kendali: \qty{2}{kg/ha} dalam \qty{8200}{m^3}: \qty{0.24}{mg/L}.
Yang ditebang: \qty{120}{kg} dalam \qty{11500}{m^3}: \qty{10.4}{mg/L},
empat puluh kali lebih banyak.
\textbf{9.} Kendali: $6.5 - 2 = +\qty{4.5}{kg/ha}$ setahun, menimbun.
Yang ditebang: $6.5 - 120 = -\qty{113.5}{kg/ha}$ setahun, kehilangan.
\textbf{10.} $120/3850 = \qty{3.1}{\%}$ dalam setahun; lubuknya akan
bertahan sekitar \num{30} tahun pada laju itu --- tetapi lajunya turun
seiring pertumbuhan kembalinya pulih.
\textbf{11.} Kalsiumnya sepuluh kali, kaliumnya delapan belas kali.
Kalium dipegang sebagian besarnya dalam \omterm{def:b1:body-plans-tissues:tissue}{jaringan} hidup dan mudah
tercuci dari \omterm{def:b1:ecosystem-dynamics:decomposition}{seresahnya}; kalsium juga terikat pada lempung tanahnya
dan dalam batunya, lalu dilepaskan lebih lambat.
\textbf{12.} Dari mineralisasi \omterm{def:b1:ecosystem-dynamics:decomposition}{seresah}, humus dan sisa tebangannya:
amonifikasi (N organik menjadi amonium) oleh \omterm{def:b1:ecosystem-dynamics:decomposition}{penguraiannya}, lalu
nitrifikasi (amonium menjadi nitrat) oleh bakteri penitrifikasi, yang
lebih cepat dalam tanah yang lebih hangat, lebih basah dan terbuka.
\textbf{13.} Turun dari \num{120} ke \num{20} ke \num{5}: di bawah
masukan \qty{6.5}{kg/ha}-nya sekitar tahun 6; serapan sebuah tutupan
baru berupa ceri pin, frambos dan anakan pohon memulihkan
penahanannya.
\textbf{14.} $\log A$: 0, 1, 2, 3, 4; $\log S$: 1.08, 1.32, 1.58,
1,83, 2,08 --- sebuah garis lurus.
\textbf{15.} $z = (2.08 - 1.08)/4 = 0.25$; $c = 12$. Pulau tengahnya:
$12\times 10^{0.25} = 21.3$, $12\times 10^{0.5} = 37.9$, $12\times
10^{0.75} = 67.5$: semuanya dalam satu spesies.
\textbf{16.} $12\times 250^{0.25} = 12\times 3.98 = 48$ spesies.
\textbf{17.} $38 = 6\times 300/(6 + 300e)$: $6 + 300e = 47.4$, $e =
\qty{0.138}{yr^{-1}}$; pergantiannya $0.138\times 38 = 5.2$ spesies setahun.
\textbf{18.} $3\times 300/(3 + 41.4) = 20$ spesies.
\textbf{19.} $30\times 300/(30 + 41.4) = 126$ spesies. Kepunahan di
sebuah pulau \qty{100}{km^2} tetap cepat: bahkan dengan jalan masuk
yang bebas, spesies yang tak sanggup memelihara sebuah \omterm{def:b1:populations:population}{populasi} pada
luas itu terus mati; luasnya, bukan jalan masuknya, yang menetapkan
langit-langitnya.
\textbf{20.} Sebuah $z$ yang lebih tinggi berarti cacah spesiesnya
turun lebih cepat bersama luasnya: pulau kecilnya relatif lebih
miskin. Secara biologis, \omterm{def:b1:populations:population}{populasi} di pulau kecil punah lebih mudah,
atau pulaunya begitu terasing sehingga sedikit spesies mencapainya
(keterasingan menaikkan $z$).
\textbf{21.} $(1000/100\,000)^{0.25} = 0.01^{0.25} = 0.32$: sekitar
sepertiga.
\textbf{22.} Ketika terasing, masing-masing memegang $(100/100\,000)^{0.25} = 0.18$
spesiesnya; bila daftarnya bebas satu sama lain keduanya akan bersama
memegang lebih banyak daripada satu cagarnya, tetapi daftarnya
bertumpang tindih berat dan masing-masing kehilangan spesies yang
memerlukan lebih dari \qty{100}{km^2}, sehingga lebih sedikit pada
kenyataannya. Ketika terhubung koridor keduanya berperilaku sebagai
satu cagar \qty{1000}{km^2} bagi spesies yang memakai koridornya:
sekitar sepertiga lagi, dengan keuntungan bahwa sebuah malapetaka
setempat tidak mengosongkan semuanya.
\textbf{23.} \qty{1000}{km^2} adalah \qty{31.6}{km} tiap sisinya:
bagian dalamnya $(31.6 - 0.6)^2 = \qty{961}{km^2}$, \qty{96}{\%}.
\qty{100}{km^2} adalah \qty{10}{km} tiap sisinya: $(10 - 0.6)^2 = 88$,
\qty{88}{\%}.
\textbf{24.} Sebuah penurunan lambat dari cacah aslinya menuju
kesetimbangan barunya $cA^z$, sepanjang berpuluh tahun: spesies yang
akan punah masih hadir pada mulanya, dalam \omterm{def:b1:populations:population}{populasi} yang terlalu kecil
untuk bertahan --- utang kepunahannya adalah beda antara cacah
sekarangnya dan kesetimbangannya, yang masih harus dibayar.
\textbf{25.} $z = 0.25$; pulau \qty{100}{km^2}-nya memegang \num{38}
spesies pada kesetimbangan dan sekitar \num{126} begitu dijembatani.
\end{solution}
